\medskip

On considère le cube ABCDEFGH ci-dessous tel que AB $= 1$.

On note M le centre de la face BCGF et N le centre de la face EFGH.

\begin{center}
\psset{unit=1cm}
\begin{pspicture}(6,5.6)
\psframe(0.3,0.3)(4,4)
\psline(4,0.3)(5.7,1.4)(5.7,5.1)(4,4)
\psline(5.7,5.1)(2,5.1)(0.3,4)
\psline[linestyle=dashed](0.3,0.3)(2,1.4)(2,5.1)
\psline[linestyle=dashed](2,1.4)(5.7,1.4)
\uput[u](2,5.1){A} \uput[u](5.7,5.1){B} \uput[r](5.7,1.4){C} \uput[dr](2,1.4){D}
\uput[l](0.3,4){E} \uput[u](4,4){F} \uput[dr](4,0.3){G} \uput[dl](0.3,0.3){H}
\end{pspicture}
\end{center}

\medskip

On se place dans le repère orthonormé $\left(\text{D}~;~ \vect{\text{DH}},~\vect{\text{DC}},~\vect{\text{DA}}\right)$.

\medskip

\begin{enumerate}
\item Donner sans justifier les coordonnées des points F et C.
\item Calculer les coordonnées des points M et N.
\item 
	\begin{enumerate}
		\item Démontrer que le vecteur $\vect{\text{AG}}$ est normal au plan (HFC).
		\item En déduire une équation cartésienne du plan (HFC).
	\end{enumerate}
\item Déterminer une représentation paramétrique de la droite (AG).
\item Démontrer que le point R de coordonnées $\left(\dfrac23~;~\dfrac23;~\dfrac13\right)$ est le projeté orthogonal du point G sur le plan (HFC).
\item On admet qu'une représentation paramétrique de la droite (FG) est:

\[\left\{\begin{array}{l c l}
x&=&1\\
y&=&1\\ 
z&=&t
\end{array}\right.\:\:(t \in \R).\]

Démontrer qu'il existe un unique point K sur la droite (FG) tel que le triangle KMN soit rectangle en K.
\item Quelle fraction du volume du cube ABCDEFGH le volume du tétraèdre FNKM représente-t-il ?
\end{enumerate}
